\scanpage{28}
% archive sheet number: 27
% L1--7; vertical emphasis stroke beside the statement
\noindent\vrule\hspace{.6em}\begin{minipage}{.93\linewidth}
\uline{Théorème} Soit $G$ un $Y$-schéma en groupes \inserted{{\unsure{commu.}}} de type fini \deleted{sur ${\unsure{Y|S}}$}, et $P^1$ un fibré principal homogène sous $G$ [devenant trivial par extension fid. plate quasi-compacte{\unsure{---}}], d’où un schéma en groupes $\mathbb Z$-augmenté $P$, avec $P^{(1)}=P^1$, $P^{(0)}=G$. Considérons l’application canonique $u{\unsure{\longrightarrow|\mapsto}}u\circ{\unsure{\varphi|\phi}}_S$
\end{minipage}\par
\[
\operatorname{Hom}_{\substack{Y\text{-schémas en groupes}\\\mathbb Z\text{-augmentés}}}
({\unsure{\mathscr P|\mathcal P}}_{S,X/Y},P)\longrightarrow\operatorname{Hom}_Y(X-S,P^1)
\]
% L8--9
\uline{Cette application est bijective}.

(Th. de Rosenlicht)

% L10
\uline{Corollaire 1}
\[
\operatorname{Hom}_{Y\text{-schémas en groupes}}(\mathscr P_{S,X/Y},G)\xrightarrow{\sim}\operatorname{Hom}_Y(X-S,G)
\]
% L11--15; vertical emphasis stroke beside lower paragraph
\uline{Corollaire 2}
\noindent\vrule\hspace{.6em}\begin{minipage}{.93\linewidth}
\inserted{Considérons} Le \deleted{{\unsure{homomorphisme de}}} \inserted{foncteur} $\varphi_S^*$, de la catégorie des schémas \uline{extensions} \inserted{(en groupes commutatifs)} de $P_{S,X/Y}$ par des $Y$-schémas en groupes commutatifs de type fini $G$ sur $Y$, dans la catégorie \deleted{{\unsure{qui}}}
\end{minipage}\par
% margin note written slanted, read after rotating clockwise
\margin{il faut faire {\unsure{attention aux définitions}} de ``fibré'' et ``extension'' ? $G$ plat sur ${\unsure{Y|S}}$??}
